% SOURCE_STATEMENT: "Can arbitrarily small perturbations of anti–de Sitter initial data, evolved by the vacuum Einstein equations with reflecting conformal-boundary conditions, form black-hole regions at sufficiently late times?"
% Preparation and final self-audit: CONTRIBUTING.md prompts C and D; source/audit artifacts retained in the local physics_statements_2026_10_11 batch.
% !TeX program = lualatex
\documentclass[11pt]{article}
\usepackage[margin=25mm]{geometry}
\usepackage{fontspec}
\setmainfont{Latin Modern Roman}
\usepackage{amsmath,amssymb,mathtools}
\usepackage{unicode-math}
\setmathfont{Latin Modern Math}
\usepackage{xurl,hyperref}
\hypersetup{colorlinks=true,urlcolor=blue}
\setcounter{secnumdepth}{0}
\setlength{\parindent}{0pt}
\setlength{\parskip}{5pt}
\setlength{\emergencystretch}{3em}
\newcommand{\sref}[2]{\hyperlink{source:#1:#2}{[S#2]}}
\newcommand{\cataloguescope}[1]{\paragraph{Recorded target.}#1}
\newcommand{\R}{\mathbb R}
\newcommand{\C}{\mathbb C}
\newcommand{\N}{\mathbb N}
\newcommand{\Z}{\mathbb Z}
\newcommand{\Q}{\mathbb Q}
\newcommand{\F}{\mathbb F}
\newcommand{\eps}{\varepsilon}
\begin{document}
\section{Nonlinear Instability of Anti–de Sitter Spacetime}\label{statement}

\subsection{Definitions and mathematical statement}

\paragraph{Shared notation.}$\mathbb N=\{1,2,\ldots\}$; the standard number fields have their usual meanings. All additional conventions are stated locally.


\paragraph{Definitions and conventions.}
Anti-de Sitter space solves the vacuum Einstein equations with negative cosmological constant, $\operatorname{Ric}(g)=\frac{2\Lambda}{D-2}g$ in spacetime dimension $D$, with $\Lambda<0$. Its conformal boundary is timelike; reflecting boundary conditions prevent energy from escaping there. Arbitrarily small perturbations means that for every positive size threshold in a specified initial-data norm one seeks smaller constraint-compatible data. A black-hole region is defined relative to causal communication with the conformal boundary.

\paragraph{Target retained from the canonical source catalogue.}
\paragraph{Catalogue wording retained for traceability.}
Can arbitrarily small perturbations of anti–de Sitter initial data, evolved by the vacuum Einstein equations with reflecting conformal-boundary conditions, form black-hole regions at sufficiently late times?

\paragraph{Source scope and qualification.}
This target is the vacuum problem. Scalar-matter examples or resonant approximations, including IDs 30002150 and 30002920, are related mechanisms rather than proofs of the stated vacuum instability. The norm and boundary regularity must be specified in a theorem.
\subsection{Short English statement}
Can sufficiently late black holes arise from vacuum anti-de Sitter perturbations of arbitrarily small size under reflecting boundary conditions?
\subsection{Sources}
\begin{itemize}
\item[\hypertarget{source:statement:1}{S1}] ULAM.AI and the UnsolvedMath Contributors, UnsolvedMath, catalogue entry. Catalogue curation under CC BY 4.0; definitions and English exposition added locally.
\url{https://huggingface.co/datasets/ulamai/UnsolvedMath}.
\item[\hypertarget{source:statement:2}{S2}] Mathematical General Relativity (2018), Oberwolfach Reports, DOI 10.4171/owr/2018/36, pp. 2157--2251. Primary report accessed 2026-10-11; relevant PDF pages 49.
\url{https://doi.org/10.4171/owr/2018/36}.
\item[\hypertarget{source:statement:3}{S3}] Bizon--Rostworowski 2011 (catalogue-listed context; not independently consulted in this preparation).
\url{https://arxiv.org/abs/1104.3702}.
\end{itemize}
\label{end:statement}
\end{document}
